\documentclass[11pt,a4paper]{article}
\usepackage{fontspec}
\setmainfont{DejaVu Sans}[Scale=0.90]
\usepackage{amsmath,amssymb}
\usepackage[margin=2.4cm]{geometry}
\usepackage{booktabs,array,longtable,tabularx}
\usepackage{enumitem}
\usepackage{titlesec}
\usepackage{parskip}
\titlelabel{\thetitle.\quad}
\titleformat{\section}{\normalfont\large\bfseries}{\thesection.}{0.6em}{}
\titleformat{\subsection}{\normalfont\normalsize\bfseries}{\thesubsection}{0.6em}{}
\titlespacing*{\section}{0pt}{16pt}{6pt}
\titlespacing*{\subsection}{0pt}{12pt}{4pt}
\setlist[itemize]{leftmargin=1.4em,itemsep=1pt,topsep=3pt}
\newcolumntype{L}[1]{>{\raggedright\arraybackslash}p{#1}}
\begin{document}
\begin{center}
{\Large\bfseries Rappels mathématiques et physiques}\\[4pt]
{\small Krapez, \emph{International Journal of Thermal Sciences} \textbf{136} (2018) 182--199}\\[2pt]
{\small Version corrigée --- 8 septembre 2026}
\end{center}

space{4pt}


\section{Grandeurs thermiques}

\subsection{Définitions}

\begin{center}
\begin{tabularx}{\textwidth}{@{}llX@{}}
\toprule
\textbf{Grandeur} & \textbf{Unité} & \textbf{Signification} \\
\midrule
Conductivité $\lambda$ & $\mathrm{W\,m^{-1}K^{-1}}$ & Flux transmis pour un gradient donné en régime
permanent \\
\addlinespace
Capacité volumique $\rho c$ & $\mathrm{J\,m^{-3}K^{-1}}$ & Énergie requise pour élever d'un kelvin
l'unité de volume \\
\addlinespace
Diffusivité $a = \lambda/\rho c$ & $\mathrm{m^{2}s^{-1}}$ & Vitesse de propagation d'une perturbation
thermique \\
\addlinespace
Effusivité $b = \sqrt{\lambda\rho c}$ & $\mathrm{W\,s^{1/2}m^{-2}K^{-1}}$ & Aptitude de la surface à
échanger de la chaleur ; gouverne l'élévation de température sous flux imposé \\
\bottomrule
\end{tabularx}
\end{center}

\subsection{Relations inverses}

\[ \lambda = b\sqrt{a} , \qquad \rho c = \frac{b}{\sqrt{a}} . \]

Ces relations se vérifient directement :
\[ b\sqrt{a} = \sqrt{\lambda\rho c}\cdot\sqrt{\lambda/\rho c} = \sqrt{\lambda^{2}} = \lambda ,
\qquad
\frac{b}{\sqrt{a}} = \sqrt{\lambda\rho c \cdot \frac{\rho c}{\lambda}} = \rho c . \]

Il en découle les deux corollaires employés dans la transformation de Liouville :
\[ \frac{\lambda}{\sqrt{a}} = b , \qquad \frac{1}{\rho c\,\sqrt{a}} = \frac{1}{b} . \]

Deux propriétés parmi les quatre suffisent à décrire complètement le milieu ; les deux autres s'en
déduisent.

\subsection{Ordres de grandeur}

\begin{center}
\begin{tabular}{@{}lcc@{}}
\toprule
\textbf{Matériau} & $a\ (\mathrm{m^{2}s^{-1}})$ & $b\ (\mathrm{W\,s^{1/2}m^{-2}K^{-1}})$ \\
\midrule
Cuivre & $1{,}1\times10^{-4}$ & $\sim 3{,}7\times10^{4}$ \\
Acier & $1{,}2\times10^{-5}$ & $\sim 1{,}4\times10^{4}$ \\
Alumine dense & $\sim 1\times10^{-5}$ & $\sim 8\times10^{3}$ \\
Zircone dense & $\sim 7\times10^{-7}$ & $\sim 3\times10^{3}$ \\
Verre & $3{,}4\times10^{-7}$ & $\sim 1{,}4\times10^{3}$ \\
Bois & $1{,}5\times10^{-7}$ & $\sim 4\times10^{2}$ \\
\bottomrule
\end{tabular}
\end{center}

Les diffusivités s'étalent sur près de trois décades entre métaux et isolants. Les valeurs numériques
précises doivent être associées aux données expérimentales de l'étude considérée.

\section{Onde thermique}

\subsection{Longueur de diffusion}

En régime périodique dans un milieu semi-infini homogène, l'amplitude décroît exponentiellement avec la
profondeur et la phase varie linéairement. La profondeur caractéristique est
\[ \mu = \sqrt{\frac{2a}{\omega}} . \]

La fréquence de modulation constitue ainsi un paramètre direct de sélection de la profondeur sondée :
les hautes fréquences ne sondent que la subsurface, les basses fréquences atteignent les couches
profondes.

\subsection{Réponse de référence}

Pour un milieu semi-infini homogène soumis à un flux modulé d'amplitude $P$,
\[ |\theta_{0}| = \frac{P}{b\sqrt{2\pi f}} , \qquad \arg(\theta_{0}) = -45^{\circ} . \]

Cette phase constante de $-45^{\circ}$ constitue la référence de la métrologie photothermique :
l'ensemble de l'analyse repose sur la mesure des écarts à cette valeur.

\subsection{Effusivité apparente}

Le contraste inverse d'amplitude par rapport à ce milieu de référence définit une effusivité apparente
dépendant de la fréquence. Elle tend vers l'effusivité de surface aux hautes fréquences et vers celle
du substrat aux basses fréquences. Elle constitue une lecture directe, mais approximative, des
variations d'effusivité en profondeur.

\section{Équations essentielles}

\subsection{Conduction unidimensionnelle}

\[ \varphi = -\lambda\,\frac{\partial T}{\partial z} ,
\qquad
\rho c\,\frac{\partial T}{\partial t} = -\frac{\partial\varphi}{\partial z} , \]
d'où, pour un milieu homogène,
\[ \frac{\partial T}{\partial t} = a\,\frac{\partial^{2}T}{\partial z^{2}} , \]
et, pour un milieu à propriétés variables après transformation temporelle,
\[ p\,\rho c(z)\,\theta = \frac{d}{dz}\!\left[\lambda(z)\,\frac{d\theta}{dz}\right] . \]

\subsection{Coordonnée et forme normale}

\[ \xi(z) = \int_{0}^{z}\frac{du}{\sqrt{a(u)}} ,
\qquad
\phi = -\,b(\xi)\,\theta' ,
\qquad
\frac{d^{2}\psi}{d\xi^{2}} - \bigl(V(\xi)+p\bigr)\psi = 0 , \]
avec $\psi = s\,\theta$, $s = b^{\pm 1/2}$ et $V = s''/s$.

\subsection{Quadripôle}

\[ \begin{pmatrix}\theta_{0}\\ \phi_{0}\end{pmatrix}
= \mathbf{M}\begin{pmatrix}\theta_{1}\\ \phi_{1}\end{pmatrix} ,
\qquad
\mathbf{M} = \begin{pmatrix} A & B \\ C & D \end{pmatrix} ,
\qquad
AD - BC = 1 . \]

Pour un empilement, $\mathbf{M}_{\text{total}}$ est le produit ordonné des matrices individuelles. La
réponse de surface d'un système revêtement-substrat s'écrit
\[ \theta_{0} = \wp\,\frac{AZ + B}{CZ + D + h\,(AZ+B)} ,
\qquad Z = \frac{1}{b_{1}\sqrt{p}} , \]
où $Z$ désigne l'impédance thermique du substrat semi-infini, $\wp$ la densité de puissance de la
source et $h$ le coefficient d'échange en face avant.

\section{Outils mathématiques}

\subsection{Équations différentielles du second ordre}

Une équation linéaire du second ordre admet exactement deux solutions indépendantes ; toute solution en
est une combinaison linéaire. Deux conditions aux limites fixent les deux constantes.

\subsection{Wronskien et théorème d'Abel}

\[ W(f,g) = f g' - f' g . \]

Le wronskien mesure l'indépendance de deux fonctions : il s'annule si et seulement si elles sont
proportionnelles. Pour une équation en forme normale, dépourvue de terme en dérivée première, le
théorème d'Abel établit que le wronskien de deux solutions est constant. Cette propriété est à
l'origine du déterminant unité de la matrice quadripolaire.

\subsection{Forme normale de Liouville}

Toute équation du type $y'' + P(x)y' + Q(x)y = 0$ se ramène à la forme $u'' + I(x)u = 0$ par la
substitution $y = u\exp\bigl(-\tfrac{1}{2}\!\int P\bigr)$. Le facteur exponentiel est déterminé par la
condition d'annulation du terme en dérivée première ; il n'est pas choisi.

Appliquée au cas thermique où $P = b'/b$, cette substitution donne $y = u\,b^{-1/2}$, c'est-à-dire
$\psi = \sqrt{b}\,\theta$.

\subsection{Transformées intégrales}

La transformée de Laplace ou de Fourier remplace la dérivation temporelle par une multiplication par
$p$, en vertu de la proportionnalité entre l'exponentielle et sa dérivée. Elle réduit ainsi l'ordre du
problème d'une variable. La condition initiale nulle est requise pour que le terme de bord de la
transformée de Laplace s'annule.

L'inversion analytique n'étant généralement pas accessible en régime transitoire, elle s'effectue
numériquement, par exemple par la méthode de De Hoog.

\section{Éléments de mécanique quantique}

\subsection{Équation de Schrödinger stationnaire}

\[ -\frac{\hbar^{2}}{2m}\,\frac{d^{2}\psi}{dx^{2}} + V(x)\,\psi = E\,\psi , \]
où $\hbar$ est la constante de Planck réduite, $V$ l'énergie potentielle et $E$ l'énergie de la
particule. Le terme en $\hbar^{2}/2m$ est l'opérateur d'énergie cinétique, obtenu en substituant
$\hat p = -i\hbar\,d/dx$ dans l'expression classique $p^{2}/2m$.

\subsection{Interprétation de la fonction d'onde}

Le module au carré de $\psi$ est une densité de probabilité de présence, d'où la condition de
normalisation $\int|\psi|^{2}dx = 1$. Cette condition d'intégrabilité de carré est à l'origine de la
quantification : seules certaines valeurs de $E$ admettent une solution acceptable, appelées valeurs
propres.

\subsection{États liés et états de diffusion}

Pour $E$ inférieure au potentiel à l'infini, la particule est confinée, la fonction d'onde décroît
exponentiellement et le spectre est discret. Pour $E$ supérieure, la fonction d'onde oscille à
l'infini, le spectre est continu et l'analyse porte sur les coefficients de réflexion et de
transmission. Le problème inverse associé, consistant à reconstruire le potentiel à partir des données
de diffusion, relève de la théorie de Gelfand, Levitan et Marchenko.

\subsection{Statut de l'analogie}

La correspondance avec la forme normale de Liouville est formelle. Dans le cas thermique, $\psi$ n'est
pas une amplitude de probabilité, aucune condition de normalisation ne s'applique, le potentiel a la
dimension d'un inverse de temps et le paramètre $E = -p$ est complexe. Les méthodes de construction de
solutions restent transposables ; les interprétations physiques ne le sont pas.
\end{document}